% Part: lambda-calculus
% Chapter: lambda-definability
% Section: pairs

\documentclass[../../../include/open-logic-section]{subfiles}

\begin{document}

\olfileid{lam}{ldf}{pai}
\olsection{جوړې او مخکنۍ شمېره}

\begin{defn}
د $M$ او $N$ جوړه (چې $\tuple{M,N}$ ليکل کېږي) داسې تعريفېږي:
\[
\tuple{M,N} \ident \lambd[f][fMN].
\]
\end{defn}
  
په شهودي ډول دا يوه تابع ده چې بله تابع اخلي او هغه
د جوړې پر دوو غړو تطبيقوي۔ د همدې نظر له مخې لاندې
جوړوونکے لرو: دوه ترمونه اخلي او د هغو لرونکې جوړه ورکوي:
\[
  \fn{Pair} \ident \lambd[mn][\lambd[f][fmn]]
\]
که يوه جوړه راکړل شي، د غړو بېرته ترلاسه کول ئې هم غواړو۔
د دې دپاره دوو لاسرسي تابعو ته اړتيا ده چې جوړه د آرګومېنټ
په توګه اخلي او لومړے يا دويم غړے ئې ورکوي:
\begin{align*}
  \fn{Fst} & \ident \lambd[p][p(\lambd[mn][m])]\\
  \fn{Snd} & \ident \lambd[p][p(\lambd[mn][n])]
\end{align*}

\begin{prob}
  څرګنده کړئ چې د لاسرسي تابعې $\fn{Fst}$ او $\fn{Snd}$
  ولې کار کوي۔
\end{prob}

اوس د جوړو په مرسته د مخکنۍ شمېرې تابع !!{lambda define}
کولے شو:
\[
  \fn{Pred} \ident \lambd[n][\fn{Fst}(n (\lambd[p][\tuple{\fn{Snd}\, {p}, \fn{Succ}(\fn{Snd}\, {p})}]) \tuple{\num 0, \num 0})]
\]
په ياد ولرئ چې $\num n\, f x$، $f^{n}(x)$ ته راکمېږي؛
دلته $f$ داسې تابع ده چې جوړه $p$ اخلي او يوه نوې جوړه
ورکوي۔ د نوې جوړې لومړے غړے د~$p$ دويم غړے او دويم
غړے ئې د هماغه دويم غړي ورپسې شمېره ده؛ $x$،
$\tuple{0,0}$ ده۔ نو که $n=0$ وي، پايله
$\tuple{0,0}$ ده، او که نۀ وي، $\tuple{\num{n-1}, \num n}$
ده۔ بيا $\fn{Pred}$ د پايلې لومړے غړے ورکوي۔

تفريق داسې تعريفولے شو چې $\fn{Pred}$ پر $a$ باندې
$b$ ځله تطبيق کړو:
\[
\fn{Sub} \ident \lambd[ab][b \fn{Pred}\, a].
\]
    
\end{document}
